\section{Introduction}

Interactive 3D anatomy viewers are an appealing format for plant biology outreach and
education: a user rotates a whole organism, clicks an organ, and reads what it is. The
format's origin, however, matters for whether what they read is true. \texttt{rice-atlas}
(\url{https://github.com/dr-richard-barker/rice-atlas}) is one such viewer, built on a
procedural functional-structural plant model and forked twice from a generic human-anatomy
template before being relabelled for rice. Its own README is explicit that the geometry is
conceptual, and inspecting the repository directly turns up no cited \textit{Oryza sativa}
dataset, no license the GitHub license detector can resolve (recorded as
\texttt{NOASSERTION}), no FAIR metadata, and no GitHub Pages deployment.

None of that makes the viewer worthless -- a clearly-labelled conceptual model is a
legitimate teaching tool. It does mean that building "the \textit{Arabidopsis} version" of
it is an underspecified request until a decision is made about which of two very different
projects is wanted: a relabelled geometric demo, or a resource that earns the adjectives
this project set out with -- FAIR, reviewed, and scientifically accurate. We chose, and
report here, a hybrid: keep the interactive organ-selectable viewer, because that is the
part of rice-atlas actually worth learning from, but rebuild its content so that every
organ's geometry and every data value shown traces to a real, checkable source.

\textit{Arabidopsis thaliana} is a fortunate choice of second species for this exercise.
Unlike rice, it has a 79-organ, whole-development RNA-seq atlas
\citep{Klepikova2016}, single-cell developmental atlases for at least two organ systems
\citep{Shahan2022,Vijayan2021}, and -- because it is the standard model organism for plant
spaceflight biology -- a substantial NASA Open Science Data Repository (OSDR) footprint,
including studies that profile root and whole-seedling transcriptomes under real and
simulated microgravity. That last point gives the resulting atlas a dimension rice-atlas
has no counterpart for at all: a spaceflight-response layer, grounded in real flight
experiments rather than invented.

We report three things: what we found rice-atlas actually to be, once inspected rather
than taken at face value; what real data sources were available to ground an
\textit{Arabidopsis} equivalent, and which ones we could and could not actually use;
and what building to a pre-registered FAIR/evidence-gate standard changed about the
result compared to simply forking the original.
